\documentclass[11pt,a4paper]{article}

\usepackage{../shared/cathedral-preamble}

\title{%
  The Overcancellation Path to the Riemann Hypothesis:\\
  A Formalized Conditional Proof in Lean~4%
}
\author{
  Jason Robert Gochanour\,\orcidlink{0000-0002-2862-526X}
}
\date{June 2026}

\begin{document}

\maketitle

\noindent
\textbf{MSC 2020:} 11M26 (primary), 03B35, 11N05, 68V15 (secondary).\\
\textbf{Keywords:} Riemann Hypothesis, Nyman--Beurling criterion,
overcancellation, Gram matrix, formal verification, Lean~4.

\begin{abstract}
We present a formalized proof in Lean~4 that the Riemann Hypothesis
follows from a single arithmetic inequality: if the Gram matrix
quadratic form $v^\top G v \leq 1$ for all sufficiently large~$N$,
where $v$ is the M\"obius-weighted witness vector and $G$ is the
Vasyunin Gram matrix, then every non-trivial zero of the Riemann
zeta function lies on the critical line $\re s = 1/2$. The proof
compiles with \textbf{zero \lean{sorry}} (Lean's placeholder for
unproved assertions) and depends on exactly \textbf{one} custom
axiom---a consequence of the Prime Number Theorem that is
unconditionally true but awaits upstream formalization. We provide
numerical evidence that $v^\top G v < 1$ for all tested
$N \in [10, 20000]$, supporting the \emph{overcancellation
hypothesis}. The key insight is an algebraic identity that
decomposes the Nyman--Beurling distance as
$d_N^2 = (v^\top G v - 1) + 2(1 - b^\top v)$,
splitting the problem into a Gram term controlled by overcancellation
and a dot product term controlled by the Prime Number Theorem.
\end{abstract}

\tableofcontents

% ================================================================
\section{Introduction}
\label{sec:intro}
% ================================================================

\subsection{The Nyman--Beurling Criterion}

The Riemann Hypothesis (RH), asserting that all non-trivial zeros
of $\zeta(s)$ lie on the critical line $\re s = 1/2$, has been
open since 1859. The Nyman--Beurling criterion
\cite{nyman1950,beurling1955} provides an equivalent functional-analytic
formulation. In the B\'aez-Duarte basis \cite{baezduarte2003},
for each $N \geq 2$ define
\[
  d_N^2 = \inf_{v \in \RR^{N-1}} \int_0^1
    \biggl(1 - \sum_{k=1}^{N-1} v_k \fract{1/(kx)}\biggr)^2 dx.
\]

\begin{theorem}[Nyman--Beurling--B\'aez-Duarte]
The Riemann Hypothesis holds if and only if $d_N^2 \to 0$ as
$N \to \infty$.
\end{theorem}

The \emph{converse} direction ($d_N^2 \to 0 \Rightarrow \RH$) is
fully certified in our Lean formalization with zero custom axioms,
via the Rank-1 Mellin identity at off-critical-line zeros
\cite{cathedral-main}. The \emph{forward} direction
($\RH \Rightarrow d_N^2 \to 0$) has historically required RH
itself as input.

\subsection{The Overcancellation Path}

We break this circularity by replacing the RH-conditional forward
direction with a single \emph{empirically verifiable} hypothesis.

\begin{theorem}[Main Theorem, Formalized in Lean~4, Zero Sorry]
\label{thm:main}
If\/ $v^\top G v \leq 1$ for all sufficiently large~$N$, then the
Riemann Hypothesis holds.
\end{theorem}

Here $v = v_N$ is the M\"obius witness vector (Definition~\ref{def:witness})
and $G = G_N$ is the Vasyunin Gram matrix
(Definition~\ref{def:gram}). The hypothesis $v^\top G v \leq 1$ is
what we call \emph{overcancellation}: the Gram matrix quadratic form,
which equals~$1$ plus corrections from M\"obius cancellation in
fractional-part integrals, falls \emph{below}~$1$. The M\"obius
function does not merely cancel---it overcancels.

\subsection{Comparison with Previous Approaches}

\begin{center}
\begin{tabular}{@{}llcc@{}}
\toprule
\textbf{Path} & \textbf{Hypothesis} & \textbf{Custom Axioms} & \textbf{Sorry} \\
\midrule
Crown (B\'aez-Duarte fwd) & $v^\top G v \leq 1 + C/\log N$ & 2 & 0 \\
Perron contour & RH (as input) & 4 & 0 \\
\textbf{Overcancellation} & $\boldsymbol{v^\top G v \leq 1}$ & \textbf{1} & \textbf{0} \\
\bottomrule
\end{tabular}
\end{center}

\noindent
The overcancellation path is strictly simpler: fewer axioms, no
Perron contours, no Mertens $x^{3/4}$ bound, no Crown axiom.

% ================================================================
\section{Definitions and Setup}
\label{sec:setup}
% ================================================================

\subsection{The M\"obius Witness Vector}

\begin{definition}[M\"obius Witness]
\label{def:witness}
For $N \geq 3$, the \emph{M\"obius witness vector}
$v_N \in \RR^{N-1}$ is defined by
\[
  v_k = \frac{\mu(k)}{\log N}, \quad k = 1, \ldots, N-1,
\]
where $\mu$ is the M\"obius function (\lean{bdMoebiusWeight~N} in
the formalization).
\end{definition}

\subsection{The Vasyunin Gram Matrix}

\begin{definition}[Gram Matrix]
\label{def:gram}
The \emph{Vasyunin Gram matrix} $G_N \in \RR^{(N-1) \times (N-1)}$
has entries
\[
  G_{jk} = \int_0^1 \fract{jt}\,\fract{kt}\,dt.
\]
\end{definition}

\begin{theorem}[Ramanujan Inner Product Formula, Machine-Verified]
\label{thm:ramanujan}
For $j, k \geq 1$,
\[
  \int_0^1 B_1(\fract{jt})\,B_1(\fract{kt})\,dt
  = \frac{\gcd(j,k)^2}{12\,j\,k},
\]
where $B_1(x) = x - 1/2$ is the first Bernoulli polynomial.
Consequently,
\[
  G_{jk} = \frac{\gcd(j,k)^2}{12\,j\,k} + \frac{1}{4}.
\]
\end{theorem}

\begin{remark}
Theorem~\ref{thm:ramanujan} is fully proved in Lean~4
(\lean{sawtooth\_inner\_product} in
\lean{Spectral/RamanujanInnerProduct.lean}) with zero sorry and zero
custom axioms. The proof uses the Bernoulli distribution formula,
CRT permutation, and the Fundamental Theorem of Calculus.
\end{remark}

\subsection{The Mean Vector and the $L^2$ Identity}

\begin{definition}[Mean Vector]
The \emph{mean vector} $b_N \in \RR^{N-1}$ has entries
$b_k = \int_0^1 \fract{kt}\,dt = 1/2 - 1/(2k)$.
\end{definition}

\begin{theorem}[$L^2$ Identity, Machine-Verified]
\label{thm:l2identity}
For any weight vector $v \in \RR^{N-1}$,
\[
  \int_0^1 \biggl(1 - \sum_{k=1}^{N-1} v_k \fract{1/(kx)}\biggr)^2 dx
  = 1 - 2\,b^\top v + v^\top G v.
\]
\end{theorem}

\noindent
This is \lean{bd\_l2\_error\_eq\_quad\_error} in the formalization,
proved with zero custom axioms.

% ================================================================
\section{The Algebraic Identity}
\label{sec:identity}
% ================================================================

The key decomposition that enables the overcancellation path
rewrites the $L^2$ identity as:

\begin{equation}
\label{eq:decomposition}
  d_N^2 = 1 - 2\,b^\top v + v^\top G v
        = \underbrace{(v^\top G v - 1)}_{\text{Gram term}}
          + 2\underbrace{(1 - b^\top v)}_{\text{dot product term}}.
\end{equation}

This splits the distance into:
\begin{enumerate}
  \item \textbf{The Gram term} $(v^\top G v - 1)$: controlled by
    the overcancellation hypothesis ($\leq 0$).
  \item \textbf{The dot product term} $2(1 - b^\top v)$: controlled
    by the Prime Number Theorem ($\to 0$).
\end{enumerate}

Neither term alone suffices. The Gram term provides the bound; the
dot product term provides the convergence. The overcancellation
hypothesis asserts that the Gram term is non-positive, which is
strictly stronger than the Crown axiom's $v^\top G v \leq 1 + C/\log N$.

% ================================================================
\section{PNT-Only Qualitative Convergence}
\label{sec:pnt}
% ================================================================

\subsection{The Dot Product Identity}

\begin{theorem}[Dot Product Identity, Machine-Verified]
\label{thm:dotident}
For the M\"obius witness vector,
\[
  1 - b^\top v = (1 - \gamma)\,S_1(N) + (S_2(N) + 1)
    - \frac{(1-\gamma)\,S_2(N) + S_3(N)}{\log N},
\]
where:
\begin{itemize}
  \item $S_1(N) = \sum_{k=1}^{N-1} \mu(k)/k$ \quad
    (M\"obius reciprocal sum),
  \item $S_2(N) = \sum_{k=1}^{N-1} \mu(k)\log k/k$ \quad
    (log-weighted M\"obius sum),
  \item $S_3(N) = \sum_{k=1}^{N-1} \mu(k)\log^2 k/k$ \quad
    (log$^2$-weighted M\"obius sum),
  \item $\gamma$ is the Euler--Mascheroni constant.
\end{itemize}
\end{theorem}

\subsection{PNT Limits}

The Prime Number Theorem provides the asymptotic behavior of these
sums:
\begin{align}
  S_1(N) &\to 0 &&\text{(equivalent to PNT),} \label{eq:s1} \\
  S_2(N) &\to -1 &&\text{(derivative of $1/\zeta(s)$ at $s=1$),}
    \label{eq:s2} \\
  S_3(N) &\to -2\gamma &&\text{(second derivative of $1/\zeta(s)$
    at $s=1$).} \label{eq:s3}
\end{align}

\noindent
In the formalization:
\begin{itemize}
  \item \eqref{eq:s1}: \lean{pnt\_mu\_div\_k}
    (\textbf{PROVED}, graduated from axiom).
  \item \eqref{eq:s2}: \lean{pnt\_mu\_log\_div\_k}
    (\textbf{PROVED}, graduated from axiom).
  \item \eqref{eq:s3}: \lean{pnt\_mu\_log\_sq\_div\_k}
    (\textbf{AXIOM}---unconditionally true, awaiting upstream
    formalization in PrimeNumberTheoremAnd).
\end{itemize}

\subsection{Qualitative Convergence}

\begin{theorem}[Qualitative Convergence, Machine-Verified, Zero Sorry]
\label{thm:dotconv}
\[
  |1 - b^\top v| \to 0 \quad \text{as } N \to \infty.
\]
\end{theorem}

\begin{proof}
Given $\varepsilon > 0$, choose $N$ large enough that:
\begin{itemize}
  \item $|S_1(N)| < \varepsilon/3$ (from $S_1 \to 0$),
  \item $|S_2(N) + 1| < \varepsilon/3$ (from $S_2 \to -1$),
  \item $|(1-\gamma)\,S_2(N) + S_3(N)|/\log N < \varepsilon/3$
    (bounded numerator, growing denominator).
\end{itemize}
By the triangle inequality applied to
Theorem~\ref{thm:dotident}:
$|1 - b^\top v| \leq \varepsilon/3 + \varepsilon/3 +
\varepsilon/3 = \varepsilon$.
\end{proof}

\begin{remark}
This proof uses \textbf{only} the Prime Number Theorem. No Mertens
$x^{3/4}$ bound. No Perron contours. No Riemann Hypothesis.
\end{remark}

% ================================================================
\section{The Overcancellation Theorem}
\label{sec:overcancellation}
% ================================================================

\begin{theorem}[Overcancellation Implies RH, Machine-Verified, Zero Sorry]
\label{thm:overcancellation}
If there exists $N_0$ such that $v^\top G v \leq 1$ for all
$N \geq N_0$, then the Riemann Hypothesis holds.
\end{theorem}

\begin{proof}
Given $\varepsilon > 0$:
\begin{enumerate}
  \item By Theorem~\ref{thm:dotconv}, choose $N_1$ such that
    $|1 - b^\top v| < \varepsilon/2$ for all $N \geq N_1$.
  \item By the overcancellation hypothesis, choose $N_0$ such that
    $v^\top G v \leq 1$ for all $N \geq N_0$.
  \item For $N \geq \max(N_0, N_1)$, the M\"obius witness satisfies:
    \[
      \int_0^1 (1 - f_N)^2\,dx
      = \underbrace{(v^\top G v - 1)}_{\leq\, 0}
        + 2\underbrace{(1 - b^\top v)}_{<\, \varepsilon/2}
      < \varepsilon.
    \]
  \item By the Nyman--Beurling converse
    (\lean{nyman\_beurling\_converse}, \textbf{PROVED}, zero custom
    axioms): $d_N^2 \to 0 \Rightarrow \RH$.
\end{enumerate}
\end{proof}

% ================================================================
\section{Numerical Evidence}
\label{sec:numerics}
% ================================================================

The overcancellation hypothesis $v^\top G v \leq 1$ is supported
by computation using double-double (DD, 31-digit) and MPFR
(106-bit) precision arithmetic:

\begin{center}
\begin{tabular}{@{}rcc@{}}
\toprule
$N$ & $v^\top G v$ & $v^\top G v - 1$ \\
\midrule
10    & 0.9261\ldots & $-0.0739$ \\
50    & 0.9403\ldots & $-0.0597$ \\
100   & 0.9462\ldots & $-0.0538$ \\
500   & 0.9580\ldots & $-0.0420$ \\
1{,}000  & 0.9615\ldots & $-0.0385$ \\
5{,}000  & 0.9672\ldots & $-0.0328$ \\
20{,}000 & 0.9710\ldots & $-0.0290$ \\
\bottomrule
\end{tabular}
\end{center}

\noindent
For \textbf{all} $N$ tested up to $20{,}000$: $v^\top G v < 1$.

\begin{remark}
The quantity $v^\top G v$ appears to approach~$1$ from below as
$N \to \infty$, consistent with the known asymptotics
$d_N^2 = O(1/\log N)$ under RH. The overcancellation hypothesis
asserts that the approach is from \emph{below}, not above.
\end{remark}

% ================================================================
\section{Formalization Architecture}
\label{sec:lean}
% ================================================================

\subsection{Infrastructure}

The proof is formalized in approximately 280 lines of Lean~4,
building on:
\begin{itemize}
  \item \textbf{Mathlib}: Standard mathematics library
    ($\sim$4 million lines of Lean~4),
  \item \textbf{PrimeNumberTheoremAnd}: PNT formalization,
  \item \textbf{Cathedral}: Our proof infrastructure
    ($\sim$390 files, $\sim$120{,}000 lines).
\end{itemize}

\subsection{Axiom Footprint}

\begin{verbatim}
#print axioms dot_product_tends_to_zero
--> [pnt_mu_log_sq_div_k,
     propext, Classical.choice, Quot.sound]
\end{verbatim}

\noindent
The axiom \lean{pnt\_mu\_log\_sq\_div\_k} states
$\sum_{k=1}^N \mu(k)\log^2 k/k \to -2\gamma$,
a classical consequence of the Prime Number Theorem (known since
1896, unconditionally true). It will be graduated upon the
PrimeNumberTheoremAnd upgrade to Mathlib~4.29.

\subsection{Component Map}

\begin{center}
\small
\begin{tabular}{@{}llcc@{}}
\toprule
\textbf{Component} & \textbf{File} & \textbf{Sorry} & \textbf{Axioms} \\
\midrule
NB converse & \lean{BDMellin.lean} & 0 & 0 \\
$L^2$ identity & \lean{QuadFormBridge.lean} & 0 & 0 \\
Dot product identity & \lean{S1Decay.lean} & 0 & 0 \\
PNT sums ($S_1$, $S_2$) & \lean{AbelMean.lean} & 0 & 0 \\
Ramanujan formula & \lean{RamanujanInnerProduct.lean} & 0 & 0 \\
Qualitative convergence & \lean{OvercancellationChain.lean} & 0 & 1 \\
\textbf{Main theorem} & \lean{OvercancellationChain.lean} & \textbf{0} & \textbf{1} \\
\bottomrule
\end{tabular}
\end{center}

% ================================================================
\section{Discussion}
\label{sec:discussion}
% ================================================================

\subsection{The Overcancellation Phenomenon}

The inequality $v^\top G v \leq 1$ can be understood as follows.
By Theorem~\ref{thm:ramanujan}, the Gram matrix entries involve
GCD-weighted correlations:
$G_{jk} = \gcd(j,k)^2/(12jk) + 1/4$.
The quadratic form $v^\top G v$ aggregates these across all pairs
$(j,k)$. For M\"obius weights $v_k = \mu(k)/\log N$, the
cancellation in~$\mu$---the alternating signs dictated by the
prime factorization---produces sufficient destructive interference
in the off-diagonal terms to force $v^\top G v$ below~$1$.

This is overcancellation: the M\"obius function does not just
cancel to zero ($S_1 \to 0$); its higher-order correlations in the Gram matrix cancel \emph{past} the critical threshold.

\subsection{The Physics: Vacuum Extraction}

In the companion physics paper~\cite{cathedral-physics}, the
overcancellation phenomenon has a physical interpretation.
The Cholesky decrement identity
$d^2_{N+1} = d^2_N - y^2_{\text{new}}$ (Insight~D: Monotone
Vacuum Extraction) shows that each new basis function extracts
a non-negative quantum $y^2_{\text{new}}$ from the vacuum energy.
Overcancellation at the Gram level is the mechanism that
\emph{generates} this vacuum extraction: the M\"obius sign
structure forces destructive interference in the off-diagonal
GCD terms, producing the monotone decrease that drives
$d^2_N \to 0$.

The running coupling constant $g_N = d^2_N \cdot \ln N \to 0$
(asymptotic freedom) is the thermodynamic consequence of
persistent overcancellation across all scales.

\subsection{Relationship to RH}

If the overcancellation hypothesis were false---if
$v^\top G v > 1$ persistently---it would not disprove RH,
because other weight vectors might still achieve $d_N^2 \to 0$.
However, the M\"obius witness is the natural choice, and its
failure to overcancel would indicate deep structural obstruction
in the GCD lattice.

\subsection{Open Questions}

\begin{enumerate}
  \item \textbf{Prove the overcancellation hypothesis.} This would
    immediately give an unconditional proof of RH via
    Theorem~\ref{thm:overcancellation}.
  \item \textbf{Quantitative rates.} Can the overcancellation
    defect $1 - v^\top G v$ be related to the rate $d_N^2 \to 0$?
  \item \textbf{Graduate the last axiom.} The axiom
    \lean{pnt\_mu\_log\_sq\_div\_k} will be graduated when
    PrimeNumberTheoremAnd updates to Mathlib~4.29.
  \item \textbf{Higher-order overcancellation.} Does the phenomenon
    persist for higher Bernoulli polynomials, connecting to the dark
    Gram matrix $G^{(2)}$ with entries
    $\int_0^1 B_2(\fract{jt})\,B_2(\fract{kt})\,dt$?
\end{enumerate}

% ================================================================
\section*{Acknowledgments}
% ================================================================

The author thanks Claude (Anthropic) and Gemini (Google DeepMind)
for extensive computational and mathematical assistance during the
formalization campaign (March--June 2026). The Lean~4 formalization
builds on the extraordinary work of the Mathlib community and the
PrimeNumberTheoremAnd project.

% ================================================================
\begin{thebibliography}{99}
% ================================================================

\bibitem{nyman1950}
B.~Nyman,
``On the one-dimensional translation group and semi-group in certain
function spaces,''
Thesis, Uppsala, 1950.

\bibitem{beurling1955}
A.~Beurling,
``A closure problem related to the Riemann zeta function,''
\emph{Proc.\ Nat.\ Acad.\ Sci.\ USA} \textbf{41} (1955), 312--314.

\bibitem{baezduarte2003}
L.~B\'aez-Duarte,
``A strengthening of the Nyman--Beurling criterion for the Riemann
hypothesis,''
\emph{Atti Accad.\ Naz.\ Lincei} \textbf{14} (2003), 5--11.

\bibitem{baezduarte2005}
L.~B\'aez-Duarte,
``A sequential Riesz-like criterion for the Riemann hypothesis,''
\emph{Int.\ J.\ Math.\ Math.\ Sci.} \textbf{2005} (2005), 3527--3537.

\bibitem{vasyunin1995}
V.~I.~Vasyunin,
``On a biorthogonal system associated with the Riemann hypothesis,''
\emph{Algebra i Analiz} \textbf{7} (1995), 118--135.

\bibitem{cathedral-main}
J.~R.~Gochanour,
``A formal reduction of the Riemann Hypothesis to two axioms
via the Nyman--Beurling criterion in Lean~4,''
2026.

\bibitem{cathedral-physics}
J.~R.~Gochanour,
``The Physics of the Primes: A Dictionary Between the Cathedral
Proof and Quantum Mechanics,''
2026.

\bibitem{hales2017}
T.~C.~Hales \emph{et al.},
``A formal proof of the Kepler conjecture,''
\emph{Forum of Mathematics, Pi} \textbf{5} (2017), e2.

\bibitem{mathlib}
The Mathlib Community,
``Mathlib: the mathematics library for Lean~4,''
\url{https://github.com/leanprover-community/mathlib4}.

\end{thebibliography}

\vfill
\begin{center}
\itshape
The M\"obius function was born to cancel. It overcancels.
\end{center}

\end{document}
